% Part: computability
% Chapter: recursive-functions
% Section: other-recursions

\documentclass[../../../include/open-logic-section]{subfiles}

\begin{document}

% SEGMENT OLP-0222-S01

\olfileid{cmp}{rec}{ore}
\olsection{பிற மீள்வரையறைகள்}

சோடியாக்கத்தையும் தொடராக்கத்தையும் பயன்படுத்தி, மேலும் வழக்கத்திற்கு
மாறானதும் பயனுள்ளதுமான முதன்மை மீள்வரையறை வடிவங்களை நியாயப்படுத்தலாம்.
எடுத்துக்காட்டாக, பின்வரும் வரையறையில் இருப்பதுபோல், இரு சார்புகளை
ஒரேநேரத்தில் வரையறுப்பது பலவேளைகளில் பயனுள்ளதாகும்:
\begin{align*}
h_0(\vec x, 0) & = f_0(\vec x) \\
h_1(\vec x, 0) & = f_1(\vec x) \\
h_0(\vec x, y+1) & = g_0(\vec x, y, h_0(\vec x, y), h_1(\vec x, y)) \\
h_1(\vec x, y+1) & = g_1(\vec x, y, h_0(\vec x, y), h_1(\vec x, y))
\end{align*}
இது \emph{ஒரேநேர மீள்வரையறையின்} ஒரு நிகழ்வு. சார்புகளை
வரையறுப்பதற்கான மற்றொரு பயனுள்ள வழி, பின்வரும் வரையறையில் உள்ளதுபோல்,
$h(\vec x, y+1)$ இன் மதிப்பை $h(\vec x, 0)$, \dots,
~$h(\vec x, y)$ ஆகிய \emph{அனைத்து} மதிப்புகளையும் கொண்டு தருவதாகும்:
\begin{align*}
h(\vec x, 0) & = f(\vec x) \\
h(\vec x, y+1) & = g(\vec x, y, \tuple{h(\vec x, 0), \dots, h(\vec x, y)}).
\end{align*}
இக்கருத்தை மேலும் சுருக்கமாகப் பின்வரும் திட்டம் பற்றிக்கொள்கிறது:
\[
h(\vec x, y) = g(\vec x, y, \tuple{h(\vec x, 0), \dots, h(\vec x, y-1)})
\]
$y$ ஆனது $0$ ஆக இருக்கும்போது $g$ இன் கடைசி உள்ளீடு வெற்றுத்
தொடர் மட்டுமே என்ற புரிதலுடன் இதைப் பயன்படுத்துகிறோம். இரு
வடிவாக்கங்களிலும், “அடுத்தெண் படி”யைக் கணிக்கும்போது, இதுவரை
கணிக்கப்பட்ட மதிப்புகளின் முழுத் தொடரையும் சார்பு $h$ பயன்படுத்தலாம்
என்பதே கருத்து. இது \emph{மதிப்புப்பாதை மீள்வரையறை} என அறியப்படுகிறது.
ஒரு குறிப்பிட்ட எடுத்துக்காட்டாக, பின்வரும் வகை வரையறையை நியாயப்படுத்த
இதைப் பயன்படுத்தலாம்:
\begin{align*}
h(\vec x, y) & = \begin{cases}
  g(\vec x, y, h(\vec x, k(\vec x, y))) & \text{$k(\vec x, y) < y$ எனில்} \\
  f(\vec x) & \text{இல்லையெனில்}
\end{cases}
\end{align*}
வேறுவிதமாகக் கூறினால்,~$k$ தரும் \emph{எந்தவொரு} முந்தைய
மதிப்பிலுள்ள $h$ இன் மதிப்பைக் கொண்டு, $y$ இலுள்ள $h$ இன்
மதிப்பைக் கணிக்கலாம்.


\begin{prob}
  மீதிச் சார்பு $r(x,y)$ ஐ மதிப்புப்பாதை மீள்வரையறையால் வரையறுக்க.
  ($x$, $y$ இயல் எண்களாகவும் $y > 0$ ஆகவும் இருந்தால்,
  ஏதேனும்~$z$ க்கு $x = z\times y + r(x,y)$ என அமையும்
  ~$y$ ஐவிடக் குறைந்த எண் $r(x,y)$. உறுதிப்பாட்டிற்காக,
  $y=0$ எனில் $r(x,0) = 0$ என்று கொள்வோம்.)
\end{prob}

வழக்கமான முதன்மை மீள்வரையறையைப் பயன்படுத்தி இச்சார்புகளை எவ்வாறு
பெறுவது என்று சிந்திக்க வேண்டும். முதன்மை மீள்வரையறையின் ஓர் இறுதி
வடிவம், செல்லும் வழியில் \emph{அளபுருக்களை} (பக்க மதிப்புகளை)
மாற்ற அனுமதிப்பதால் மேலும் நெகிழ்வானது:
\begin{align*}
h(\vec x, 0) & = f(\vec x) \\
h(\vec x, y+1) & = g(\vec x, y, h(k(\vec x), y))
\end{align*}
இதையும் வழக்கமான முதன்மை மீள்வரையறையால் உருவகப்படுத்தலாம்.
(அவ்வாறு செய்வது சிக்கலானது. ஒரு குறிப்புக்கு, கணக்கீட்டை கையால்
பின்னோக்கி விரித்துப் பார்க்க.)

\end{document}
